\section{Evidence distance does not predict responses to planted errors}
\label{sec:reversal}

% source: PAPER_SKELETON Sec 3; EVIDENCE_protocol_prereg P5.
% NOTE (r2 framing directive #1 + #2): governance vocabulary (pre-registration,
% Holm family, H1'/H2'/H4 code names, the scoreboard figure, evidence-status
% labels) removed from this section and consolidated in Appendix~\ref{app:provenance};
% directive #2 recasts the framing paragraphs from research-journey voice into
% flat findings about the models. Numbers unchanged.

\paragraph{A distance gradient does not predict rejection or reuse.} A natural
hypothesis is that errors get caught when the disproving evidence is nearby:
rejection and reuse should both be graded by how far the refuting evidence
sits from the plant. A matched-distance test on Qwen2.5-7B, holding everything
but distance fixed across six distance-linked predictions, finds no support
for it. Rejection is flat to slightly rising across distance rather than
falling (slope $+0.011$, $p{=}0.79$), and reuse is highest one hop away rather
than farthest (slope $-0.248$, $p{=}0.996$); both run opposite to a gradient.
% source: EVIDENCE_protocol_prereg P5 (H1' slope +0.011 p=0.79; H4 slope -0.248 p=0.996)
The one distance-linked prediction the test supports is not a gradient but a
polarity control: affirmative and negated false plants are absorbed
equivalently, within a $\pm0.10$ margin (diff $+0.054$, $p{=}0.0012$).
% source: EVIDENCE_protocol_prereg P5 (H2' PASS p=0.0012, diff +0.054); PAPER_SKELETON 3.1

\paragraph{The polarity equivalence holds under the stricter validator too.} A
reader might worry that this surviving result rests on the permissive
closure-valid standard we elsewhere demote. It does not. Under the strict
hop-sound standard the polarity equivalence still holds within the
$\pm0.10$ margin, with a difference of $-0.042$ and a TOST
$p{=}8\times10^{-5}$; the point sign flips but the equivalence conclusion is
unchanged.
% source: PAPER_SKELETON 3.2; EVIDENCE_protocol_prereg P5 (hop-sound H2' diff -0.042, TOST p=8e-5)
We state this as ``equivalent within $\pm0.10$'', never as ``identical'',
because the equivalence fails at the tighter $\pm0.05$ margin and the base
rates are low.
% source: EVIDENCE_protocol_prereg P5 caveat (fails at 0.05; aff 0.108 / neg 0.153 hop-sound)

\paragraph{Two downstream contrasts carry different evidential weight.} Two
contrasts the rest of the paper builds on differ in how directly the
matched-distance test speaks to them, and we mark the difference rather than
treating them alike. The $d{=}0$ visibility contrast of
Section~\ref{sec:visibility} is measured directly but was left
direction-open, so it motivates the visibility picture rather than clinching
it.
% source: EVIDENCE_protocol_prereg P5 caveat (C-0 non-family); PAPER_SKELETON 3.3
The reversed reuse trend of Section~\ref{sec:usability}, where reuse is
highest one hop away, runs opposite to the gradient prediction and does not
reach a confirmatory threshold under matched construction, so we report it as
a suggestive observation and never describe it as a confirmed effect.
% source: EVIDENCE_protocol_prereg P5 caveat (H4 reversal one-sided p=0.004, exploratory)

\paragraph{No strict metric reverses a conclusion the paper relies on.} We
scanned the contrasts the paper leans on under the strict hop-sound
standard and found no relied-on reversal; the one sign flip, in the
polarity equivalence, strengthens rather than weakens the case, because it
moves the estimate toward zero difference.
% source: PAPER_SKELETON 3.4; EVIDENCE_protocol_prereg P4 (reversal scan, 5 contrasts)
The strict tables are in Appendix~\ref{app:strict}.
